% ===================== INTRODUCCIÓN =====================
\section{Introducción}

La industria azucarera colombiana genera anualmente entre 6 y 7 millones de toneladas de bagazo de caña, de las cuales 5 a 6 millones se destinan a la cogeneración de energía eléctrica y vapor de proceso \cite{sac2020}. Las calderas bagaceras acuotubulares constituyen el equipo central de estos sistemas de cogeneración.

El presente informe desarrolla el cálculo termodinámico detallado para determinar el consumo de bagazo requerido por una caldera acuotubular que debe producir 100~t/h de vapor sobrecalentado a \SI{106}{barg} y \SI{545}{\celsius}. Se emplean datos experimentales de la composición elemental del bagazo colombiano reportados en la literatura científica, en lugar de valores genéricos de libros de texto, con el fin de obtener resultados más representativos de las condiciones reales de operación de los ingenios azucareros del Valle del Cauca.
